\documentclass[11pt]{article}
\usepackage{fontspec,amsmath,amssymb,amsthm,hyperref}
\usepackage[margin=1in]{geometry}
\setlength{\emergencystretch}{3em}
\title{Finite entropy supplies the gradient moment directly}
\author{LangevinLean mathematical exposition}
\date{}
\begin{document}\maketitle
\noindent Source module: \texttt{UnadjustedLangevin/LSI/EntropyMoment.lean}. Compile this proof document with XeLaTeX or LuaLaTeX. The accompanying HTML pairs the mathematical argument with the exact Lean statements and records the scope of the formalization.


A finite initial KL divergence already controls the squared Gibbs score when \(\nabla U\) is Lipschitz. Gaussian averaging of the directional score MGF gives \(\mathbb E_\nu e^{\|\nabla U\|^2/(4L)}\le e^{d/2}\). The entropy inequality then proves both integrability and \(G_\mu\le4L H(\mu\mid\nu)+2dL\). This direct argument closes the moment input to ULA without assuming an entropy-cost coupling or a transport inequality.


\section{mul le klFun add exp sub one}


Let \(\Phi(a)=a\log a-a+1\), with its continuous value \(\Phi(0)=1\). For \(a>0\), the elementary inequality \(e^u\ge1+u\), applied to \(u=b-\log a\), gives \(e^b\ge a+ab-a\log a\). Rearranging proves \(ab\le\Phi(a)+e^b-1\). The zero-density case is handled separately: its left side is zero, and its right side is \(e^b\ge0\). Thus this entropy/exponential Young inequality remains valid at zero Radon-Nikodym density.



\noindent Lean declarations: \texttt{mul\_le\_klFun\_add\_exp\_sub\_one}.

\section{integrable of finite kl of integrable exp}


Write \(a=d\mu/d\nu\). Finite KL supplies absolute continuity \(\mu\ll\nu\) and integrability of the actual log likelihood; equivalently \(\Phi(a)\) is integrable under \(\nu\). For a measurable nonnegative observable \(f\), the Young inequality gives \(0\le af\le\Phi(a)+e^f\). Both terms on the right are integrable by the finite-KL and exponential-moment hypotheses.



Domination proves \(af\in L^1(\nu)\). The Radon-Nikodym weighted-integrability equivalence then gives \(f\in L^1(\mu)\). In particular, the moment conclusion is derived before any ordinary expectation is manipulated; it is not an extra assumption hidden in the entropy inequality.



\noindent Lean declarations: \texttt{integrable\_of\_finite\_kl\_of\_integrable\_exp}.

\section{integral le kl add of integral exp le}


Assume \(\int e^f d\nu\le e^c\). First derive \(f\in L^1(\mu)\) with the preceding theorem. Apply the Young inequality to \(a=d\mu/d\nu\) and \(f-c\), and integrate: \(\int(f-c)d\mu\le\int\Phi(a)d\nu+\int e^{f-c}d\nu-1\).



Probability normalization identifies \(\int\Phi(a)d\nu\) with the genuine finite KL divergence and \(\int(f-c)d\mu=\int f d\mu-c\). The MGF hypothesis makes \(\int e^{f-c}d\nu\le1\), canceling the final two terms. Thus \(\int f d\mu\le H(\mu\mid\nu)+c\), together with the integrability witness used in its derivation. The proof explicitly justifies each weighted integral and subtraction.



\noindent Lean declarations: \texttt{integral\_le\_kl\_add\_of\_integral\_exp\_le}.

\section{integrable radial gaussian}


For \(b>0\), the real radial Gaussian \(e^{-b\|v\|^2}\) is integrable on every finite-dimensional real Hilbert space. The proof specializes the complex Gaussian Fourier integrability theorem to real positive precision and takes its norm; the complex norm is exactly the displayed real exponential. This supplies the finite integral needed in the Gaussian linearization argument.



\noindent Lean declarations: \texttt{integrable\_radial\_gaussian}.

\section{lintegral radial gaussian}


Because the radial Gaussian is nonnegative and integrable, its nonnegative Lebesgue integral equals the embedded ordinary integral. The Gaussian Fourier formula evaluates the latter as \((\pi/b)^{d/2}\). Hence \(\int^- e^{-b\|v\|^2}dv=\operatorname{ofReal}((\pi/b)^{d/2})\), a positive finite normalization factor that can safely be canceled later.



\noindent Lean declarations: \texttt{lintegral\_radial\_gaussian}.

\section{lintegral radial gaussian mul exp inner}


Complete the square in \(-b\|v\|^2+\langle g,v\rangle\): it equals \(\|g\|^2/(4b)-b\|v-g/(2b)\|^2\). Exponentiating separates the constant factor \(e^{\|g\|^2/(4b)}\) from a translated radial Gaussian. Translation invariance of Lebesgue measure and the radial normalization give \(\int^- e^{-b\|v\|^2}e^{\langle g,v\rangle}dv=e^{\|g\|^2/(4b)}(\pi/b)^{d/2}\). The identity is proved for nonnegative integrals before any integrability of a random score is assumed.



\noindent Lean declarations: \texttt{lintegral\_radial\_gaussian\_mul\_exp\_inner}.

\section{lintegral exp norm sq le of directional mgf}


Suppose \(\mathbb E e^{\langle g,v\rangle}\le e^{L\|v\|^2/2}\) for every direction \(v\). Multiply by \(e^{-L\|v\|^2}\) and integrate in \(v\). Tonelli applies to the nonnegative jointly measurable integrand, so the order of integration can be exchanged without assuming the desired exponential moment in advance.



Gaussian linearization identifies the left side with \((\pi/L)^{d/2}\mathbb E e^{\|g\|^2/(4L)}\). The directional MGF bound controls the right side by \(\int e^{-(L/2)\|v\|^2}dv=(2\pi/L)^{d/2}\). Cancel the positive finite factor \((\pi/L)^{d/2}\) in the extended nonnegative reals to obtain \(\mathbb E e^{\|g\|^2/(4L)}\le2^{d/2}\).



\noindent Lean declarations: \texttt{lintegral\_exp\_norm\_sq\_le\_of\_directional\_mgf}.

\section{gibbs score exponential moment}


The smooth Gibbs score has the previously proved directional MGF bound \(\mathbb E_\nu e^{\langle\nabla U,v\rangle}\le e^{L\|v\|^2/2}\), obtained from translation invariance and first-order smoothness. Apply Gaussian averaging with \(g=\nabla U\). This gives a finite nonnegative exponential integral bounded by \(2^{d/2}\); measurability and this finiteness prove ordinary integrability.



Convert the nonnegative integral back to its real expectation. Since \(\log2\le1\), \(2^{d/2}\le e^{d/2}\), giving the convenient bound \(\mathbb E_\nu e^{\|\nabla U\|^2/(4L)}\le e^{d/2}\). Only differentiability, Lipschitz gradient, positive \(L\), and finite Gibbs normalization are used; neither convexity nor LSI enters this exponential-moment theorem.



\noindent Lean declarations: \texttt{gibbs\_score\_exponential\_moment}.

\section{gibbs gradient moment bound of finite kl}


Apply the entropy inequality with \(f(x)=\|\nabla U(x)\|^2/(4L)\) and \(c=d/2\). The Gibbs exponential-moment theorem provides both exponential integrability and the MGF bound; finite initial KL provides the absolute-continuity and entropy inputs. Therefore \(\mathbb E_\mu[\|\nabla U\|^2/(4L)]\le H(\mu\mid\nu)+d/2\), with this observable integrable under \(\mu\).



Multiply by \(4L>0\), proving \(G_\mu\le4L H(\mu\mid\nu)+2dL\) and actual square integrability of the score. The proof also derives finite Gibbs partition from its assumed probability normalization. This replaces the previously missing LSI-to-transport route by a stronger direct entropy/score estimate.



\noindent Lean declarations: \texttt{gibbs\_gradient\_moment\_bound\_of\_finite\_kl}.

\section{gibbs gradient moment bound of finite kl of alpha le}


If \(0<\alpha\le L\), then \(4L\le4L^2/\alpha\). Finite KL is nonnegative, so replacing the direct estimate's coefficient \(4L\) by \(4L^2/\alpha\) preserves the inequality. Thus \(G_\mu\le(4L^2/\alpha)H(\mu\mid\nu)+2dL\), the coefficient used in the ULA one-step calculation. Integrability is inherited from the stronger direct theorem. Compatibility \(\alpha\le L\) is proved from LSI in a separate module.



\noindent Lean declarations: \texttt{gibbs\_gradient\_moment\_bound\_of\_finite\_kl\_of\_alpha\_le}.

\end{document}
